\subsection{Subalgebras of the algebra of continuous functions on a compact space}

Let X be a compact topological space and B a unital subalgebra of \(\mathcal{C}(\mathrm{X})\) . We denote by ev the map \(x \mapsto ev_{x}\) from X into \(\mathrm{X}(\mathrm{B})\) such that \(\mathrm{ev}_{x}(f) = f(x)\) for all \(f \in B\) . We denote by j the injection of B into \(\mathcal{C}(\mathrm{X})\) , so \(\mathrm{X}(j) = \mathrm{ev}\) .

PROPOSITION 1. — Let \(f \mapsto \|f\|_{\mathrm{B}}\) be a norm endowing B with the structure of a Banach algebra.

\begin{enumerate}
\def\labelenumi{\alph{enumi})}
\item
  The map j has norm ≤ 1 in \(\mathcal{L}(B; \mathcal{C}(\mathrm{X}))\) ;
\item
  The radical of the algebra B is zero;
\item
  The map ev is continuous from X into X(B);
\item
  If B separates the points of X, then the map ev is a homeomorphism from X onto a closed subset of X(B).
\end{enumerate}

One can identify X with \(\mathsf{X}(\mathcal{C}(\mathsf{X}))\) and \(\mathcal{G}_{\mathcal{C}(\mathsf{X})}\) with the identity map (example 1 of I, p.~36). The map ev is then identified with \(\mathsf{X}(j)\) , which proves c).

For every function \(f \in \mathrm{B}\) and every \(x \in \mathrm{X}\) , one has \(\mathcal{G}_{\mathrm{B}}(f)(\mathrm{ev}_x) = f(x)\) , whence

\[
\| f \| _ {\mathcal {C} (\mathrm{X})} = \sup _ {x \in \mathrm{X}} | \mathcal {G} _ {\mathrm{B}} (f) (\mathrm{ev} _ {x}) | \leqslant \varrho_ {\mathrm{B}} (f) \leqslant \| f \| _ {\mathrm{B}},
\]

(cf. prop. 7 of I, p. 38, a)) which entails \(a\) ). Moreover, this shows that the Gelfand transform of B is injective, whence \(b\) ) (prop. 8 of I, p.~38). If B separates the points of X, then the map ev is injective, whence \(d\) ) since X is compact.

We now turn to the surjectivity of the map ev.

PROPOSITION 2. --- Let \(f \mapsto \|f\|_{\mathrm{B}}\) be a norm such that \((\mathrm{B}, \| \cdot \|_{\mathrm{B}})\) is a Banach algebra.

\begin{enumerate}
\def\labelenumi{\alph{enumi})}
\item
  If the map ev is surjective, then the algebra B is a full subalgebra of \(\mathcal{C}(\mathrm{X})\) ;
\item
  Suppose that B is a full subalgebra of \(\mathcal{C}(\mathrm{X})\) , and that there exists an element \(a \in \mathrm{B}\) such that the set of elements \(f(a)\) , where \(f\) ranges over the set of rational functions on \(\mathbf{C}\) having no pole on \(\mathrm{Sp}_{\mathrm{B}}(a)\) , is dense in B. Then the map ev is surjective;
\item
  If B is a full involutive subalgebra of \(\mathcal{C}(\mathrm{X})\) then the map ev is surjective.
\end{enumerate}

Assertion \(a\) ) follows from prop. 10 of I, p.~40, and assertion \(b\) ) of prop. 11 of I, p.~41, since the full subalgebra of B generated by \(a\) is the set of elements \(f(a)\) , where \(f\) ranges over the set of rational functions on \(\mathbf{C}\) having no pole on \(\mathrm{Sp}_{\mathrm{B}}(a)\) (lemma I, p.~6 and proposition 6 of I, p.~37, \(b\) )).

We prove \(c\) ). Suppose then that B is a full involutive subalgebra of \(\mathcal{C}(\mathrm{X})\) . To prove that ev is surjective, it suffices to prove that for every \(\chi \in \mathrm{X}(\mathrm{B})\) , there exists \(y \in \mathrm{X}\) such that \(\operatorname{Ker}(\chi) = \operatorname{Ker}(\mathrm{ev}_y)\) (th. 2 of I, p.~30). Let \(\mathrm{I} = \operatorname{Ker}(\chi)\) . This is a maximal ideal of B. Let \(\Phi\) be the set of \(x \in \mathrm{X}\) such that \(f(x) = 0\) for all \(f \in \mathrm{I}\) . We show that \(\Phi\) is not empty. In the contrary case, since X is compact, there would exist an integer \(n \geqslant 1\) , an open covering \((\mathrm{V}_1, \ldots, \mathrm{V}_n)\) of X and, for every integer \(i\) such that \(1 \leqslant i \leqslant n\) , a function \(f_i \in \mathrm{I}\) such that \(f_i(x) \neq 0\) for all \(x \in \mathrm{V}_i\) . Since the algebra B is an involutive subalgebra of \(\mathcal{C}(\mathrm{X})\) , the function

\[
f = \sum_ {i = 1} ^ {n} f _ {i} \overline {{f}} _ {i}
\]

would belong to I. But \(f(x) > 0\) for all \(x \in X\) , and therefore f would be invertible in \(\mathcal{C}(X)\) . Since B is assumed to be a full subalgebra of \(\mathcal{C}(X)\) , the function \(f \in I\) would be invertible in B, which is impossible. Consequently, the set \(\Phi\) is not empty. Let y be an element of \(\Phi\) ; the kernel of the character \(ev_{y}\) contains I, and is therefore equal to I.

Example. --- Let X = \([0,1]\) and let \(n \geqslant 0\) be an integer. Let B be the algebra of functions \(f: X \to C\) admitting continuous derivatives on \([0,1]\) up to order n, equipped with the norm considered in Example 4 of I, p.~18. Then B is a full involutive subalgebra of \(\mathcal{C}(X)\) separating the points of X, hence \(\mathsf{X}(\mathsf{B})\) is identified with X.

We consider the logarithm function as defined on \(R_{+}\) and taking values in \(R \cup \{-\infty\}\) by setting \(\log(0) = -\infty\) .

PROPOSITION 3. --- Let X be a compact space. Let B be a unital Banach subalgebra of \(\mathcal{C}(\mathrm{X})\) , equipped with the induced norm, separating the points of X. One identifies X with a closed subset of \(\mathrm{X(B)}\) (prop. 1, d)).

\begin{enumerate}
\def\labelenumi{\alph{enumi})}
\item
  For every \(f \in \mathbf{B}\) , the Gelfand transform \(\mathcal{G}_{\mathrm{B}}(f)\) is a continuous function in \(\mathsf{X}(\mathbf{B})\) which extends \(f\) and satisfies \(\| f \| = \sup |\mathcal{G}_{\mathrm{B}}(f)|\) . In particular, \(\mathcal{G}_{\mathrm{B}}\) is an isometric isomorphism from B onto a Banach subalgebra of \(\mathcal{C}(\mathsf{X}(\mathbf{B}))\) ;
\item
  Let B* be the set of invertible elements of B. For every character \(\chi \in \mathsf{X}(\mathsf{B})\) there exists a positive measure \(\mu\) of mass 1 on X such that, for every \(f \in \mathsf{B}^*\) , one has
\end{enumerate}

\[
\log (| \chi (f) |) = \int_ {\mathrm{X}} \log (| f |) d \mu .
\]

Moreover, for every \(f \in \mathrm{B}\) , one has

\[
\chi (f) = \int_ {\mathrm{X}} f d \mu ;
\]

\begin{enumerate}
\def\labelenumi{\alph{enumi})}
\setcounter{enumi}{2}
\tightlist
\item
  We assume that every element of \(\mathcal{C}_{\mathbf{R}}(\mathrm{X})\) is a uniform limit of real parts of functions belonging to B. Then, for every \(\chi \in \mathrm{X}(\mathrm{B})\) , there exists a unique measure \(\mu_{\chi} \geqslant 0\) on X such that, for every function \(f \in B\) , one has
\end{enumerate}

\[
\chi (f) = \int_ {\mathrm{X}} f d \mu_ {\chi}.
\]

Moreover, for every function \(f \in \mathrm{B}\) , one has

\[
\log (| \chi (f) |) \leqslant \int_ {\mathrm{X}} \log (| f |) d \mu_ {\chi};
\]

the function \(\log|f|\) being bounded above, the integral exists in \(R \cup \{-\infty\}\) .

Assertion \(a\) ) follows from the identification of X with a closed subset of \(\mathsf{X}(\mathsf{B})\) and the inequalities

\[
\| f \| = \sup _ {x \in \mathrm{X}} | f (x) | \leqslant \sup _ {\chi \in \mathrm{X} (\mathrm{B})} | \chi (f) | = \sup | \mathcal {G} _ {\mathrm{B}} (f) | = \varrho_ {\mathrm{B}} (f) \leqslant \| f \|
\]

for all \(f \in B\) .

Let \(\chi \in \mathsf{X}(\mathbf{B})\) and let \(n\) be a positive integer. Let \(\lambda_1, \ldots, \lambda_n \in \mathbf{R}\) and \(f_1, \ldots, f_n \in \mathbf{B}^*\) . We then have

\[
\sum_ {i = 1} ^ {n} \lambda_ {i} \log (| \chi (f _ {i}) |) \leqslant \sup _ {x \in X} \left(\sum_ {i = 1} ^ {n} \lambda_ {i} \log (| f _ {i} (x) |)\right).\tag{1}
\]

Indeed, by continuity, it suffices to prove this inequality when the real numbers \(\lambda_{i}\) are rational. By reducing to a common denominator, we reduce to the case where \(\lambda_{i} \in Z\) for all i. The inequality then becomes

\[
\log \left(\left| \chi \left(f _ {1} ^ {\lambda_ {1}} \dots f _ {n} ^ {\lambda_ {n}}\right) \right|\right) \leqslant \sup _ {x \in X} \log \left(\left| \left(f _ {1} ^ {\lambda_ {1}} \dots f _ {n} ^ {\lambda_ {n}}\right) (x) \right|\right),
\]

and follows from the fact that \(\|\chi\|=1\) (th. 1 of I, p.~29).

Let B' be the vector subspace of \(\mathcal{C}_{\mathbf{R}}(\mathrm{X})\) generated by the functions \(\log(|f|)\) for \(f \in B^{*}\) . The estimate (1) proves that there exists a linear form h of norm \(\leqslant 1\) on B' such that \(\log(|\chi(f)|) = h(\log(|f|))\) for all \(f \in B^{*}\) . By the Hahn--Banach theorem (EVT, II, p.~24, cor. 2), the linear form h extends to a linear form \(\mu\) of norm \(\leqslant 1\) on \(\mathcal{C}_{\mathbf{R}}(\mathrm{X})\) , that is, to a real measure \(\mu\) on X such that \(\|\mu\| \leqslant 1\) (INT, III, §1, no. 5). Taking as the element f of B* the constant \(e = \exp(1)\) we see that \(1 = \mu(1)\) . Hence, writing \(\mu = \mu^{+} - \mu^{-}\) as the difference of two mutually singular positive measures (INT, III, §1, no. 6, th. 3), it follows

\[
1 = \mu^ {+} (1) - \mu^ {-} (1) \leqslant \mu^ {+} (1) + \mu^ {-} (1) = \| \mu \| \leqslant 1,
\]

whence \(\mu = \mu^{+} \geqslant 0\) and \(\|\mu\| = 1\) .

For every \(f \in \mathrm{B}\) , one has \(\exp(f) \in \mathrm{B}^*\) , hence

\[
\begin{array}{r l} \int_ {\mathrm{X}} \mathcal {R} (f) d \mu = \int_ {\mathrm{X}} \log (| \exp (f) |) d \mu = \log (| \chi (\exp (f)) |) \\ & = \log (| \exp (\chi (f)) |) = \mathcal {R} (\chi (f)), \end{array}
\]

where we used Cor. 1 of I, p.~66. Applying this equality to \(if\) , we conclude that \(\int_{\mathrm{X}} f d\mu = \chi(f)\) . This establishes \(b\) ).

Let us place ourselves under the assumptions of \(c\) ). The existence of \(\mu_{\chi}\) follows from \(b\) . On the other hand, we have \(\mu_{\chi}(\mathcal{R}(f)) = \mathcal{R}(\chi(f))\) for all \(f \in B\) . Since the real parts of functions \(f \in B\) are dense in \(\mathcal{C}_{\mathbf{R}}(\mathrm{X})\) by hypothesis, the measure \(\mu_{\chi}\) is uniquely determined by \(\chi\) .

Let \(f \in \mathrm{B}\) . Let \(\varepsilon > 0\) be a real number. By hypothesis, there exists a function \(g \in \mathrm{B}\) such that

\[
\mathcal {R} (g) - \varepsilon \leqslant \log (| f | + \varepsilon) \leqslant \mathcal {R} (g) + \varepsilon .\tag{2}
\]

Let \(h = \exp(g) \in \mathrm{B}^*\) . By (2), we have

\[
| h | e ^ {- \varepsilon} \leqslant | f | + \varepsilon \leqslant | h | e ^ {\varepsilon}.\tag{3}
\]

The bound implies \(|fh^{-1}| \leqslant e^{\varepsilon}\) , whence \(|\chi (fh^{-1})| \leqslant e^{\varepsilon}\) and consequently

\[
\log (| \chi (f) |) \leqslant \log (| \chi (h) |) + \varepsilon = \int_ {\mathrm{X}} \log (| h |) d \mu_ {\chi} + \varepsilon .\tag{4}
\]

The lower bound in (3) then implies

\[
\log (| \chi (f) |) \leqslant \int_ {\mathrm{X}} \log (| f | + \varepsilon) d \mu_ {\chi} + 2 \varepsilon .
\]

Letting \(\varepsilon\) tend to 0, we deduce that

\[
\log (| \chi (f) |) \leqslant \int_ {\mathrm{X}} \log (| f |) d \mu_ {\chi}.
\]

This completes the proof.

